\documentclass[11pt]{article}
\usepackage[T1]{fontenc}
\usepackage[utf8]{inputenc}
\usepackage{lmodern}
\usepackage[margin=1in]{geometry}
\usepackage{amsmath,amssymb,amsthm,mathtools}
\usepackage{microtype}
\usepackage{hyperref}
\hypersetup{colorlinks=true,linkcolor=blue,citecolor=blue,urlcolor=blue,
 pdftitle={Walsh restricted isometries under sampling with replacement}}
\newtheorem{theorem}{Theorem}[section]
\newtheorem{proposition}[theorem]{Proposition}
\newtheorem{lemma}[theorem]{Lemma}
\newtheorem{corollary}[theorem]{Corollary}
\theoremstyle{remark}
\newtheorem{remark}[theorem]{Remark}
\newcommand{\E}{\mathbb E}
\newcommand{\Prob}{\mathbb P}
\newcommand{\F}{\mathbb F}
\newcommand{\R}{\mathbb R}
\newcommand{\1}{\mathbf 1}
\newcommand{\K}{\mathcal K}
\newcommand{\supp}{\operatorname{supp}}
\newcommand{\KL}{D_{\mathrm{KL}}}
\newcommand{\gaussbinom}[2]{\genfrac{[}{]}{0pt}{}{#1}{#2}_{\!2}}
\DeclareMathOperator{\Var}{Var}
\allowdisplaybreaks
\setlength{\emergencystretch}{2em}
\title{Walsh restricted isometries under sampling with replacement\\[4pt]
\large An extension of the entropic sampling argument}
\author{}
\date{Working manuscript, 30 September 2026}

\begin{document}
\maketitle

\begin{abstract}
Let $N=2^d$, and sample $m$ rows of the Walsh matrix independently and
uniformly, retaining repetitions. We determine, up to universal constants,
the smallest $m$ for which the normalized sample has restricted isometry
constant at most $1/2$ at sparsity $k$ with probability at least $0.9$:
the answer is one for $k=1$, and
$k\log(2k)\log(2eN/k)$ for $2\leq k\leq N$.
The upper bound uses the entropic encoding estimates of Burstein,
Iosevich, and Krause. We extend their sampling argument by applying the
decoder to indexed row occurrences and randomizing paired assignments to
two independent samples. The lower bound follows the Walsh subspace
obstruction of B{\l}asiok, Lopatto, Luh, Marcinek, and Rao. Averaging over
affine offsets under the fixed-sample law gives an elementary second-moment
bound uniform over all dyadic sparsities. A reduction completes the dense
and nondyadic regimes. The entropic encoding result is an explicitly
identified input from the public literature; the sampling extension and
the uniform lower-bound calculation are proved here.
\end{abstract}

\section{Statement and attribution}

All logarithms are natural unless a base is displayed. Let $G=\F_2^d$,
$d\geq1$, and $N=2^d$. Write
\[
 H(a,b)=(-1)^{a\cdot b},\qquad a,b\in G,
 \qquad H^TH=NI_N.
\]
Thus $N^{-1/2}H$ is the normalized Walsh matrix. Let
$a_1,\ldots,a_m$ be independent uniform elements of $G$ and set
\begin{equation}\label{eq:sample}
 \Phi_m=\frac1{\sqrt m}
 \begin{pmatrix}H(a_1,:)\\ \vdots\\H(a_m,:)\end{pmatrix}.
\end{equation}
Every draw is retained, including repetitions. Denote by $E_{m,N,k}$ the
event that
\begin{equation}\label{eq:rip}
 \frac12\|x\|_2^2\leq\|\Phi_m x\|_2^2\leq\frac32\|x\|_2^2
 \quad\text{for all }x\in\R^N\text{ with }|\supp x|\leq k,
\end{equation}
and define
\[
 m_*(N,k)=\min\{m\in\mathbb N:m\geq1,\ \Prob(E_{m,N,k})\geq0.9\}.
\]
This is the convention in Problem FR-10 of the public
\emph{Open Problems in Numerical Linear Algebra} catalog \cite{FR10}.

\begin{theorem}\label{thm:main}
For every $N=2^d$, $d\geq1$, one has $m_*(N,1)=1$. There are universal
constants $c,C>0$ such that, simultaneously for $2\leq k\leq N$,
\begin{equation}\label{eq:main}
 c\,k\log(2k)\log\frac{2eN}{k}
 \leq m_*(N,k)\leq
 C\,k\log(2k)\log\frac{2eN}{k}.
\end{equation}
One may take $c=1/2000$. The upper bound uses the entropic encoding
estimates of \cite[Sections 2--4]{BIK}, in the occurrence form stated
in Lemma~\ref{lem:encoding} below.
\end{theorem}

\paragraph{Attribution and scope.}
The essential improvement in the upper logarithmic rate belongs to
Burstein--Iosevich--Krause \cite{BIK}. Their stated sampling theorems
concern Bernoulli selection and uniform subsets of distinct rows. Here we
supply the passage to independent sampling with replacement, including
$m>N$. This is an extension of their proof, not an independent method
for obtaining the improved rate.

The lower-bound mechanism---a missed Walsh subspace produces a sparse
kernel vector, followed by a second-moment argument---is due to
B{\l}asiok--Lopatto--Luh--Marcinek--Rao \cite{BLLMR}. The calculation here
uses affine offsets and the exact joint law for a fixed number of
independent draws to obtain one uniform, nonasymptotic bound. We make no
claim that the underlying obstruction is new. The catalog also records
Colbrook's sharper leading-constant analyses at small fixed sparsities
and near full sparsity \cite{FR10}; the order estimate
\eqref{eq:main} does not supersede those constants.

The sole specialized external proof input is Lemma~\ref{lem:encoding}.
All other estimates needed below are proved in this manuscript. The
references are public; no unpublished companion note or private file is
a dependency. In particular, the lower bound is independent of
\cite{BIK}. Statements about the upper bound carry exactly the dependence
on that paper identified above. This working manuscript does not assert
an exhaustive priority determination for the extensions.

\section{The entropic input and its occurrence interpretation}

For this section let $H$ be any real $N\times N$ matrix with
$H^TH=NI_N$ and $|H(a,b)|\leq1$. Fix $2\leq k\leq N$ and put
\begin{equation}\label{eq:class}
 \K_k=\{y\in\R^N:\|y\|_2=1,\ \|y\|_1\leq\sqrt k\},
 \qquad L=\log\frac{2eN}{k}.
\end{equation}
For $z_a(y)=|(Hy)_a|^2$, orthogonality and the norm constraint give
\begin{equation}\label{eq:energyfacts}
 \E_a z_a(y)=1,\qquad 0\leq z_a(y)\leq k.
\end{equation}
Every $k$-sparse unit vector belongs to $\K_k$.

Choose $0<\eta<1/16$. Let $a_j=(1+\eta)^j$, $b_j=a_j^2$, and choose
integers $j_-<j_+$ satisfying
\[
 \frac{\eta}{4(1+\eta)}<a_{j_-}\leq\frac\eta4,
 \qquad
 \frac{\sqrt k}{1-\eta}\leq a_{j_+}
 <\frac{(1+\eta)\sqrt k}{1-\eta}.
\]
Write $\mathcal J=\{j_-,\ldots,j_+\}$ and $\ell=|\mathcal J|$.
Elementary geometric summation gives
\begin{align}
 \ell&\leq C_\eta\log(2k),&
 B:=\sum_j b_j&\leq C_\eta k,\label{eq:geometry1}\\
 \sum_{j\leq h}b_j&\leq C_\eta b_h,&
 \sum_j b_j(j_+-j+1)&\leq C_\eta k.\label{eq:geometry2}
\end{align}
Constants $C_\eta$ may depend only on $\eta$; constants denoted $C_0$
below are absolute.

\begin{lemma}[Entropic encoding input; Burstein--Iosevich--Krause]
\label{lem:encoding}
Fix an ordered list of $M$ rows $h_1,\ldots,h_M$ of $H$, with repetitions
allowed. For every $y\in\K_k$, there are correction records with $d_j$
positions and one of three labels at each position at level $j$, such that
\begin{equation}\label{eq:cost}
 0\leq d_j\leq M,\qquad D:=\sum_j b_jd_j\leq C_\eta kL.
\end{equation}
A deterministic decoder produces disjoint subsets $I_j\subseteq[M]$ and
$v_i=\sum_j b_j\1_{I_j}(i)$ satisfying
\begin{equation}\label{eq:approx}
 \big|v_i-|h_i y|^2\big|
 \leq C_0\eta |h_i y|^2+C_0\eta^2.
\end{equation}
The set $I_j$ is determined by the records at levels $h\geq j$.
At each level, the number of possible records of length $d$ is at most
\begin{equation}\label{eq:recordnumber}
 N_M(d)=\binom Md3^d,\qquad 0\leq d\leq M.
\end{equation}
The decoder can also be run on arbitrary such records, whether or not
they arise from a target $y$.
\end{lemma}

\begin{proof}[Justification of the occurrence interpretation]
The estimates and decoder are the public input from
\cite[Lemmas 2.2 and 3.1, Proposition 3.2, and Lemmas 4.1--4.2 and
6.1]{BIK}.
Their deterministic procedure uses an ordered sequence of bounded row
queries. Replace each query index by its occurrence index. The experts
remain $\{\pm\sqrt{k}e_b:b\in[N]\}$, so the initial comparison entropy
remains bounded by $L$, independent of $M$. Each recorded correction at
level $j$ decreases the potential by at least $3\eta^2b_j/(32k)$.
Telescoping proves \eqref{eq:cost} for the repeated-query sequence.
The decoder follows the same query order and updates, so its amplitude
estimate and the higher-level dependence identified in
\cite[Lemma 6.1]{BIK} are unchanged. Record positions
are now chosen from $[M]$, giving \eqref{eq:recordnumber}. These steps
use neither distinctness of queries nor $M\leq N$.
\end{proof}

The distinction between $M$ and $N$ will be important: $M$ counts
observed occurrences, whereas $N$ counts coordinates in the target.
No enlargement of the ambient coordinate space is involved.

\section{From entropic records to two independent samples}

The estimates in this section implement the occurrence extension of the
counting and symmetrization argument of \cite[Sections 5--8]{BIK}.
We give the probability calculation explicitly because signs attached
to paired occurrences need not be independent across individual
occurrences.

\subsection{Weighted counts}

Set $h_M(d)=\log N_M(d)$, with $h_M(0)=0$. If
$D=\sum_j b_jd_j>0$, the binomial estimate and concavity of the logarithm
give
\begin{align}
 \sum_j b_jh_M(d_j)
 &\leq\sum_{d_j>0}b_jd_j\log\frac{3eM}{d_j}\notag\\
 &\leq D\log\frac{3eM B}{D}.\label{eq:weightedentropy}
\end{align}
Indeed, apply Jensen with weights $b_jd_j/D$; the corresponding weighted
average of $1/d_j$ is at most $B/D$. If $D=0$, the left side is zero.
Since $D\leq MB$ and $u\mapsto u\log(3eMB/u)$ is increasing on
$(0,MB]$, \eqref{eq:cost} and \eqref{eq:geometry1} imply, for $M=2m$,
\begin{equation}\label{eq:entropyfinal}
 \sum_j b_jh_M(d_j)
 \leq C_\eta kL\log\left(e+\frac mL\right).
\end{equation}
For completeness, let $D_0=C_\eta kL$ be the bound in
\eqref{eq:cost}. If $D_0\leq MB$, substitute $D_0$ in
\eqref{eq:weightedentropy}; the factor $MB/D_0$ is at most a constant
depending on $\eta$ times $m/L$. If $D_0>MB$, use
$D\log(3eMB/D)\leq MB\log(3e)\leq D_0\log(3e)$.
This proves \eqref{eq:entropyfinal} in both cases.

\subsection{Paired signs and simultaneous decoded bands}

Fix an ordered pool
\[
 \Omega=\{(t,0),(t,1):1\leq t\leq m\}
\]
of $M=2m$ row occurrences. A row value may appear more than once.
Independently of this fixed pool, let $\varepsilon_1,\ldots,
\varepsilon_m$ be independent uniform signs. For a deterministic subset
$I\subseteq\Omega$, put
\[
 S_I=\sum_{t=1}^m\varepsilon_t
 \big(\1_I(t,0)-\1_I(t,1)\big).
\]
Writing the coefficient in parentheses as $c_t$, one has
\begin{equation}\label{eq:varianceproxy}
 \sum_t c_t^2\leq |I|,
 \qquad
 \E_\varepsilon e^{uS_I}
 =\prod_t\cosh(uc_t)\leq e^{u^2|I|/2}.
\end{equation}
Here $\log\cosh x\leq x^2/2$, as follows by integrating
$\tanh x\leq x$ for $x\geq0$ and using symmetry. Chernoff's argument
therefore yields, for $q>0$,
\begin{equation}\label{eq:bandtail}
 \Prob_\varepsilon\{|S_I|>\sqrt{2|I|q}\}\leq 2e^{-q}.
\end{equation}
For an empty set the asserted bound holds deterministically.

\begin{lemma}\label{lem:uniformsigned}
Fix $0<\rho<1$. With conditional probability at least $1-\rho$ over the
paired signs, simultaneously for all targets $y\in\K_k$ and the records
of Lemma~\ref{lem:encoding},
\begin{equation}\label{eq:signedv}
 \left|\sum_t\varepsilon_t(v_{t,0}-v_{t,1})\right|
 \leq\sqrt{2V A_m},\qquad V=\sum_{i\in\Omega}v_i,
\end{equation}
where one may take
\begin{equation}\label{eq:Am}
 A_m=C_\eta\left[
 kL\log\left(e+\frac mL\right)+k\log\frac{2\ell}{\rho}
 \right].
\end{equation}
The constant can be chosen independently of the pool and of $m,N,k$.
\end{lemma}

\begin{proof}
Let
\[
 Z_M=\sum_{d=0}^M N_M(d)^{-1}\leq\sum_{d=0}^M3^{-d}\leq\frac32,
 \qquad \pi_M(d)=\frac1{Z_M N_M(d)}.
\]
For a fixed $j$ and tail profile
$\mathbf d=(d_j,\ldots,d_{j_+})$, the decoded sets $I_j$ constitute a
family $\mathcal G_{j,\mathbf d}$ of cardinality at most
\[
 \prod_{h\geq j}N_M(d_h).
\]
This family is fixed before the signs are drawn. Allocate failure
probability
\[
 p_{j,\mathbf d}=\frac\rho\ell\prod_{h\geq j}\pi_M(d_h)
\]
to it. Summing these probabilities over all profiles and levels gives
$\rho$. For nonempty families, a union bound using
\eqref{eq:bandtail} shows that all their bands satisfy
$|S_I|\leq\sqrt{2|I|Q_{j,\mathbf d}}$, outside an event of probability
at most $\rho$, where
\begin{align}
 Q_{j,\mathbf d}
 &=\log\frac{2|\mathcal G_{j,\mathbf d}|}{p_{j,\mathbf d}}\notag\\
 &\leq \log\frac{2\ell}{\rho}
 +2\sum_{h\geq j}h_M(d_h)
 +(j_+-j+1)\log Z_M.\label{eq:profileQ}
\end{align}
Empty families require no bound.

Now take the records of any target $y$. By disjointness of its decoded
bands and Cauchy--Schwarz,
\begin{align*}
 \left|\sum_t\varepsilon_t(v_{t,0}-v_{t,1})\right|
 &\leq\sum_j b_j\sqrt{2|I_j|Q_{j,\mathbf d}}\\
 &\leq\sqrt{2\left(\sum_jb_j|I_j|\right)
                   \left(\sum_jb_jQ_{j,\mathbf d}\right)}.
\end{align*}
The first parenthesis equals $V$. For the second, reverse the order in
the double sum of \eqref{eq:profileQ} and apply
\eqref{eq:geometry1}--\eqref{eq:geometry2} to obtain
\[
 \sum_jb_jQ_{j,\mathbf d}
 \leq C_\eta\left[
 k\log\frac{2\ell}{\rho}+k+\sum_hb_hh_M(d_h)\right]
 \leq A_m.
\]
The last step is \eqref{eq:entropyfinal}, with the constant enlarged.
The union bound covered all records, so this event is uniform in $y$,
even if the target is selected after the signs have been observed.
\end{proof}

\subsection{The independent-copy inequality}

Let $a_{t,0},a_{t,1}$, $1\leq t\leq m$, be independent uniform row
indices. Apply the preceding lemma to this ordered pool, and use the
independent signs to assign one member of each pair to the first sample
and the other to the second. The two resulting samples are independent
lists of $m$ uniform rows. For $y\in\K_k$, let their energies be
$X_y$ and $X'_y$. Regardless of the assignments,
\[
 X_y+X'_y=\sum_{t,s} z_{a_{t,s}}(y),
 \qquad
 X_y-X'_y=\sum_t\varepsilon_t
 \big(z_{a_{t,0}}(y)-z_{a_{t,1}}(y)\big).
\]
By \eqref{eq:approx},
\[
 V\leq(1+C_0\eta)(X_y+X'_y)+2C_0\eta^2m.
\]
The difference between the signed sums of $v$ and $z$ is at most
$C_0\eta(X_y+X'_y)+2C_0\eta^2m$. Lemma~\ref{lem:uniformsigned}
therefore proves the following statement after averaging over the pool.

\begin{proposition}\label{prop:pair}
For two independent samples of $m$ uniform rows, with probability at
least $1-\rho$,
\begin{equation}\label{eq:pair}
 \begin{split}
 |X_y-X'_y|\leq{}& C_1\sqrt{A_m(X_y+X'_y+\eta^2m)}\\
 &+C_1\eta(X_y+X'_y)+C_1\eta^2m
 \qquad (y\in\K_k),
 \end{split}
\end{equation}
where $C_1$ is absolute and $A_m$ is as in \eqref{eq:Am}.
There is no restriction $m\leq N$.
\end{proposition}

The decoder's occurrence order must be fixed before the signs. In
particular, it is not obtained by first listing whichever occurrences
the signs assign to the first sample. Keeping the original paired order
is what makes the families in the union bound deterministic.

\section{The with-replacement upper bound}

\begin{lemma}[Absorption]\label{lem:absorb}
Fix $0<\delta<1/4$. There exist positive constants $c_0,c_1$, depending
only on the absolute constant $C_1$ in \eqref{eq:pair}, such that the
following holds. Suppose $u,v\geq0$, $m>0$,
$|v-m|\leq\delta m/4$, $\eta\leq c_0\delta$, $A\leq c_1\delta^2m$,
and
\[
 |u-v|\leq C_1\sqrt{A(u+v+\eta^2m)}
                +C_1\eta(u+v)+C_1\eta^2m.
\]
Then $|u-m|\leq\delta m$.
\end{lemma}

\begin{proof}
Put $s=u+v$. Since $s\leq2v+|u-v|$ and $v\leq(1+\delta/4)m$, Young's
inequality bounds the square-root term by $s/4+C_2A+C_2\eta m$, for an
absolute $C_2$. Choose $\eta$ small enough that $C_1\eta\leq1/4$.
Using $A\leq m$ then gives $s\leq C_3m$, with $C_3$ absolute.
Substitute this bound in the assumed inequality to get
\[
 |u-v|\leq C_4\sqrt{Am}+C_4\eta m.
\]
Choose $c_0,c_1$ small enough that the right side is at most
$3\delta m/4$. The triangle inequality proves the result.
\end{proof}

\begin{theorem}\label{thm:upper}
For every fixed $0<\delta<1/4$, there is $C_\delta>0$ such that
\begin{equation}\label{eq:upperm}
 m\geq C_\delta k\log\frac{2eN}{k}\log(2k)
\end{equation}
implies, with probability at least $23/24$,
\[
 \left|\frac1m\sum_{t=1}^m |(Hy)_{a_t}|^2-1\right|\leq\delta
 \qquad\text{for every }y\in\K_k.
\]
This holds for any real $H$ as above and independent uniform row draws.
\end{theorem}

\begin{proof}
Set $\rho=1/32$ and first fix $\eta>0$ small enough for
Lemma~\ref{lem:absorb}. We claim that a sufficiently large $C_\delta$ in
\eqref{eq:upperm} ensures $A_m\leq c_1\delta^2m$.
Indeed, put $t=m/(kL)$ and $q=\log(2k)$. By
\eqref{eq:geometry1} and \eqref{eq:Am},
\begin{equation}\label{eq:Aoverm}
 \frac{A_m}{m}\leq C_\eta
 \left[\frac{\log(e+kt)}t+\frac{\log(64\ell)}{Lt}\right].
\end{equation}
Here $L\geq\log(2e)$, $\ell\leq C_\eta q$, and
$\log(e+kt)\leq q+\log(e+t)$. If $t\geq Bq$, the right side of
\eqref{eq:Aoverm} tends uniformly to zero as $B\to\infty$:
$q/t\leq1/B$, $\log(e+t)/t$ is decreasing, and
$\log(64C_\eta q)/q$ is bounded for $q\geq\log4$.
This proves the claim.

For a fixed $y\in\K_k$, the independent-copy energy satisfies, by
\eqref{eq:energyfacts},
\[
 \E X'_y=m,\qquad \Var(X'_y)\leq mk.
\]
Consequently Chebyshev's inequality gives
\begin{equation}\label{eq:ghost}
 \Prob\{|X'_y-m|>\delta m/4\}
 \leq\frac{16k}{\delta^2m}\leq\frac14
\end{equation}
once $m\geq64k/\delta^2$, also ensured by increasing $C_\delta$.

Let $B$ be the event that the first sample fails the asserted uniform
estimate. On $B$, choose a witness $y$ with $|X_y-m|>\delta m$.
There is no measurability problem: the first sample takes only finitely
many values, and a witness can be chosen for each bad value. Conditional
on the first sample, \eqref{eq:ghost} applies to this fixed witness.
Thus with conditional probability at least $3/4$, its independent-copy
energy lies within $\delta m/4$ of $m$. On that event,
Lemma~\ref{lem:absorb} implies that \eqref{eq:pair} cannot hold for this
witness. Proposition~\ref{prop:pair} therefore gives
\[
 \frac34\Prob(B)\leq\rho,\qquad
 \Prob(B)\leq\frac{4\rho}{3}=\frac1{24}.
\]
This proves the theorem, with no upper restriction on the number of draws.
\end{proof}

\section{A uniform affine-subspace lower bound}

Return to the Walsh matrix on $G=\F_2^d$. Fix
\begin{equation}\label{eq:dyadicpars}
 K=2^r,\qquad 1\leq r\leq d-1,\qquad s=d-r.
\end{equation}
For an $r$-dimensional subspace $V\subseteq G$ and a linear functional
$f\in V^*$, define
\[
 A(V,f)=\{a\in G:a\cdot v=f(v)\text{ for all }v\in V\}.
\]
This affine row subspace has density $1/K$.

\begin{lemma}[A missed affine subspace]\label{lem:witness}
If none of the $m$ sampled rows lies in $A(V,f)$, then $\Phi_m$ has a
$K$-sparse unit vector in its kernel.
\end{lemma}

\begin{proof}
Define $x_b=K^{-1/2}(-1)^{f(b)}$ for $b\in V$ and $x_b=0$ otherwise.
Then $\|x\|_2=1$. Orthogonality of characters on $V$ gives
\[
 (Hx)_a=K^{-1/2}\sum_{b\in V}(-1)^{a\cdot b+f(b)}
        =\sqrt K\,\1_{A(V,f)}(a).
\]
The sampled coordinates of this vector vanish under the stated condition.
\end{proof}

Let $X$ count all pairs $(V,f)$ for which $A(V,f)$ is missed. The number
of such pairs is $T=K\gaussbinom dr$, where the bracket denotes a Gaussian
binomial coefficient. In the fixed-$m$ sampling model,
\begin{equation}\label{eq:firstmoment}
 p=(1-1/K)^m,\qquad \E X=Tp>0.
\end{equation}

\begin{lemma}[Exact offset average]\label{lem:joint}
Fix $r$-dimensional subspaces $V,W$ with $\dim(V\cap W)=j$.
Average over independent uniform $f\in V^*$ and $g\in W^*$.
The probability that both $A(V,f)$ and $A(W,g)$ are missed is
\begin{equation}\label{eq:joint}
 2^{-j}\left(1-\frac2K+\frac{2^j}{K^2}\right)^m
 +(1-2^{-j})\left(1-\frac2K\right)^m.
\end{equation}
For $j=0$, this is exactly $p^2$.
\end{lemma}

\begin{proof}
The restrictions of $f$ and $g$ to $V\cap W$ agree with probability
$2^{-j}$. In that case their common affine solution set has codimension
$\dim(V+W)=2r-j$, hence density $2^j/K^2$; otherwise it is empty.
Each of the two affine sets has density $1/K$. A single draw misses
their union with the probabilities appearing in \eqref{eq:joint}.
Independence of the $m$ draws gives the powers. At $j=0$ the first base
equals $(1-1/K)^2$, and the second term has coefficient zero.
\end{proof}

\begin{lemma}[Intersection probability]\label{lem:intersection}
Let $V,W$ be independent uniform $r$-dimensional subspaces of $\F_2^d$,
and write $q_j=\Prob\{\dim(V\cap W)=j\}$. For $1\leq j\leq r$,
\begin{equation}\label{eq:qbound}
 q_j\leq4\,2^{-j(s-r+j)}.
\end{equation}
Every feasible $j$ also satisfies $j\geq r-s$.
\end{lemma}

\begin{proof}
Fix $V$. If $\dim(V\cap W)\geq j$, then some $j$-dimensional subspace
$U\subseteq V$ lies in $W$. For a fixed $U$, elementary subspace counting
gives
\[
 \Prob\{U\subseteq W\}
 =\prod_{i=0}^{j-1}\frac{2^r-2^i}{2^d-2^i}
 \leq2^{-js}.
\]
For example, the product identity follows by counting ordered independent
$j$-tuples in a random $W$ and in the whole space. A union bound over the
$\gaussbinom rj$ possible $U$ yields
\[
 q_j\leq\gaussbinom rj 2^{-js}
 \leq4\,2^{j(r-j)}2^{-js}.
\]
For the last inequality, use
\[
 \gaussbinom rj
 =2^{j(r-j)}
 \prod_{i=0}^{j-1}\frac{1-2^{i-r}}{1-2^{i-j}}
 \leq \frac{2^{j(r-j)}}{\prod_{h=1}^\infty(1-2^{-h})}
 <4\,2^{j(r-j)}.
\]
Indeed the infinite product is at least
$(1-1/2)(1-1/4)(1-\sum_{h\geq3}2^{-h})=9/32>1/4$;
the inequality for the tail follows first for finite products and then
by a limit. Finally $\dim(V+W)\leq d$ implies $j\geq2r-d=r-s$.
\end{proof}

\begin{proposition}[Uniform second moment]\label{prop:second}
Under \eqref{eq:dyadicpars}, if $m$ is a positive integer satisfying
\begin{equation}\label{eq:lowerthreshold}
 m\leq\frac{Krs\log2}{100},
\end{equation}
then
\begin{equation}\label{eq:ratio}
 \frac{\E X^2}{(\E X)^2}
 \leq1+\frac4{2^{3/4}-1}<7.
\end{equation}
In particular, with probability greater than $1/7$, $\Phi_m$ has a
$K$-sparse unit vector in its kernel.
\end{proposition}

\begin{proof}
Using Lemma~\ref{lem:joint} to average over all ordered pairs of affine
subspaces, divide by $p^2$ from \eqref{eq:firstmoment}. Put
\[
 A_j=1+\frac{2^j-1}{(K-1)^2},\qquad
 B_0=1-\frac1{(K-1)^2}\in[0,1].
\]
The exact ratio is
\begin{equation}\label{eq:exactratio}
 R:=\frac{\E X^2}{(\E X)^2}
 =\sum_{j=0}^r q_j\left[2^{-j}A_j^m+(1-2^{-j})B_0^m\right].
\end{equation}
Since $A_0=1$ and $B_0^m\leq1$,
\begin{align}
 R&\leq1+\sum_{j\geq1}q_j2^{-j}(A_j^m-1)\notag\\
 &\leq1+\sum_{\substack{j\geq1\\q_j>0}}
 q_j2^{-j}\exp\left(\frac{m2^j}{(K-1)^2}\right).
 \label{eq:ratioexp}
\end{align}
Here and below $j\leq r$ in sums over feasible intersections.

Write $l=r-j$. Feasibility gives $0\leq l\leq s$ and $j\geq1$.
The elementary inequality
\begin{equation}\label{eq:keyelementary}
 rs\,2^{-l}\leq2j(s-l+1)
\end{equation}
holds throughout this range. For $l=0$ it is immediate. For $l\geq1$,
\[
 \frac{s}{s-l+1}\leq l,\qquad \frac{r}{j}=1+\frac lj\leq l+1,
\]
so the ratio of the left side of \eqref{eq:keyelementary} to
$j(s-l+1)$ is at most $l(l+1)2^{-l}\leq2$. The last inequality holds
at $l=1,2$ and propagates for $l\geq2$ because
$(l+2)/(2l)\leq1$.

As $K\geq2$, one has $K/(K-1)^2\leq4/K$. Combining this with
\eqref{eq:lowerthreshold} and \eqref{eq:keyelementary} yields
\begin{align}
 \frac{m2^j}{(K-1)^2}
 &\leq\frac{4\log2}{100}\,rs\,2^{j-r}\notag\\
 &\leq\frac{8\log2}{100}\,j(s-r+j+1)\notag\\
 &\leq\frac{\log2}{4}\,j(s-r+j+1).\label{eq:exponent}
\end{align}
Lemma~\ref{lem:intersection} now gives, for every term of
\eqref{eq:ratioexp},
\[
 q_j2^{-j}\exp\left(\frac{m2^j}{(K-1)^2}\right)
 \leq4\,2^{-\frac34 j(s-r+j+1)}.
\]
Every feasible $j$ has $s-r+j+1\geq1$. Summing a geometric series proves
\[
 R\leq1+4\sum_{j\geq1}2^{-3j/4}
   =1+\frac4{2^{3/4}-1}<7.
\]
The final strict inequality follows from $2^{3/4}>5/3$, which can be
checked by raising both positive sides to the fourth power.
Finally, Cauchy--Schwarz applied to $X\1_{\{X>0\}}$ gives
\[
 \Prob\{X>0\}\geq\frac{(\E X)^2}{\E X^2}>\frac17.
\]
Lemma~\ref{lem:witness} proves the kernel assertion.
\end{proof}

\begin{remark}
The offset average in \eqref{eq:joint} is taken under the actual
independent-draw law. No Poissonization or independent occupation
approximation is used. In particular, $j=0$ contributes exactly one to
the normalized pair probability. For highly overlapping subspaces,
\eqref{eq:keyelementary} controls the excess without a restriction that
$r$ or $d-r$ grow with $d$. This is why a constant failure probability
suffices uniformly up to the boundary of the dyadic range.
\end{remark}

\section{All sparsities and completion of the theorem}

\begin{proof}[Proof of Theorem~\ref{thm:main}]
Each column of $\Phi_m$ has squared norm exactly one for every $m\geq1$.
Thus $m_*(N,1)=1$.

For $2\leq k\leq N$, Theorem~\ref{thm:upper} with $\delta=1/8$ gives
\eqref{eq:rip} with probability at least $23/24>0.9$ whenever
$m\geq Ck\log(2k)\log(2eN/k)$. Taking a ceiling and enlarging the
universal constant proves the upper bound in \eqref{eq:main}.

For the lower bound suppose first that $N\geq4$. Choose
\[
 r=\min\{\lfloor\log_2 k\rfloor,d-1\},\qquad K=2^r.
\]
Then $k/2\leq K\leq k$ and $2\leq K\leq N/2$. Proposition~\ref{prop:second}
excludes every positive integer
\[
 m\leq T_0:=\frac{K\log K\log(N/K)}{100\log2}
\]
from the defining set for $m_*(N,k)$: a $K$-sparse kernel vector is also
$k$-sparse, and its probability exceeds $0.1$. Consequently
$m_*(N,k)>T_0$. If $T_0<1$, this conclusion follows instead directly
from $m_*(N,k)\geq1$.

Let $F=k\log(2k)\log(2eN/k)$. Since $k\leq2K$, $K\geq2$ and
$N/K\geq2$,
\[
 \log(2k)\leq3\log K,\qquad
 \log\frac{2eN}{k}\leq\log\frac{2eN}{K}\leq4\log\frac NK.
\]
Thus $F\leq24K\log K\log(N/K)$ and
\[
 T_0\geq\frac F{2400\log2}>\frac F{2000}.
\]
For the remaining case $N=2$, $k=2$, the inequality $m_*(2,2)\geq1$
is more than sufficient for the same lower constant. This completes
the proof for every parameter pair.
\end{proof}

\begin{remark}[Full sparsity and multiplicities]
At $k=N$, the formula gives order $N\log N$. In fact, if $n_a$ denotes
the number of times row $a$ is sampled, orthogonality diagonalizes
$\Phi_m^T\Phi_m$ with eigenvalues $Nn_a/m$. Hence the full-sparsity
event requires every occupancy to lie between $m/(2N)$ and $3m/(2N)$.
Keeping all $N$ distinct rows would give an exact isometry, but it is a
different sampling rule. The proof above retains the prescribed rule
also when $m>N$.
\end{remark}

\begin{remark}[What the extension establishes]
The result is an order estimate at fixed distortion and fixed success
probability. It does not determine the optimal upper constant, sharp
distortion dependence, or the leading constants in endpoint asymptotics.
Its upper-bound mechanism is the entropic encoding of \cite{BIK}.
The additional work is the probability interface for repeated
occurrences and the uniform affine second-moment completion of the
Walsh lower bound. Those are the contributions claimed here.
\end{remark}

\begin{thebibliography}{9}

\bibitem{BIK}
W.~Burstein, A.~Iosevich, and B.~Krause,
\emph{Restricted isometry of sampled Fourier and Hadamard matrices via
entropic descent}, arXiv:2609.22568v1 (2026).
\href{https://arxiv.org/abs/2609.22568v1}{Public preprint};
\href{https://arxiv.org/html/2609.22568v1}{versioned full text}.

\bibitem{BLLMR}
J.~B{\l}asiok, P.~Lopatto, K.~Luh, J.~Marcinek, and S.~Rao,
\emph{An improved lower bound for sparse reconstruction from subsampled
Walsh matrices}, Discrete Analysis 2023:3;
arXiv:1903.12135v3.
\href{https://arxiv.org/abs/1903.12135v3}{Public paper}.

\bibitem{FR10}
\emph{Open Problems in Numerical Linear Algebra},
FR-10: Sharp sampling complexity for Walsh restricted isometries.
Public problem statement and endpoint-results record, consulted
30 September 2026.
\href{https://github.com/ajt60gaibb/OpenProblemsInNLA/blob/main/frames-and-matrix-designs/FR-10/README.md}{Problem and bibliography}.

\end{thebibliography}
\end{document}
